\section{Operators, tensor tests, and conventions}\label{sec:conventions}

\subsection{The descendant theory}

All varieties are smooth and projective over \(\C\), unless they occur as fibers of a specified semistable family. Cohomology has complex coefficients and its usual parity. Tensor products, symmetric powers, permutations, and derivatives are taken in super vector spaces. For a target \(X\) of complex dimension \(n\), write
\[
 (a,b)_X=\int_X a\cup b,\qquad
 \mu_X|_{H^{p,q}(X)}=(p-n/2)\id,\qquad
 \mathcal R_X(a)=c_1(TX)\cup a.
\]
The pairing is perfect and even. Its coevaluation is the categorical inverse of this pairing; it includes the signs dictated by the super symmetry. We never insert a second parity involution into a node kernel.

Choose a Hodge-homogeneous basis \(\{\phi_a\}\) and let
\(t(z)=\sum_{m\geq0,a}t_m^a\phi_a z^m\), where \(t_m^a\) has the parity of \(\phi_a\). The connected potentials and total descendant potential are
\begin{align}
 F_g^X(t)
 &=\sum_{\beta}\sum_{r\geq0}
   \frac{\mathsf Q^\beta}{r!}
   \int_{[\overline{\cM}_{g,r}(X,\beta)]^\vir}
        \prod_{i=1}^r\left(\sum_{m,a}t_m^a\psi_i^m\ev_i^*\phi_a\right),
 \label{eq:potential}\\
 Z_X(t)&=\exp\left(\sum_{g\geq0}\h^{g-1}F_g^X(t)\right).
 \label{eq:partition}
\end{align}
Only stable maps contribute. The \(\psi_i\) are cotangent classes at markings on the source of the stable map, not classes pulled back by stabilization to curves. In particular, the same convention is used in relative theories and in the fiber calculation below.

Equation \eqref{eq:partition} packages disconnected domains, with the empty domain contributing \(1\). We use labelled coefficient extraction to discuss multilinear tensors: distinct formal variables distinguish all prescribed insertions, even when their values coincide. This convention accounts automatically for the factorials in \eqref{eq:potential}.

All statements are coefficientwise. Fixing an ample integral curve degree, a finite list of markings, and a genus coefficient leaves finite sums of effective curve classes and discrete stable-map data. In intermediate relative theories we impose the additional contact and end-degree bounds specified in the cap construction. Laurent series in \(\h\) are always bounded below at the finite auxiliary-variable and degree orders in use. For several replicas, the genus and curve variables are independent. The cap construction proves the corresponding boundedness when the counting polarization is only relatively ample.

\subsection{The positive Virasoro operators}

Let
\[
 \mathcal H_X=H^*(X)((z^{-1})),\qquad
 \Omega_X(f,g)=\operatorname*{Res}_{z=0}(f(-z),g(z))_X\,dz.
\]
With the standard polarization
\(\mathcal H_X=H^*(X)[z]\oplus z^{-1}H^*(X)[[z^{-1}]]\), set
\begin{equation}
 \ell_0=z\partial_z+\tfrac12+\mu_X+\frac{\mathcal R_X}{z},
 \qquad
 \ell_k=\ell_0(z\ell_0)^k\quad(k\geq1).
 \label{eq:loop-operators}
\end{equation}
These are infinitesimal symplectic operators. For example,
\(\mu_X^*=-\mu_X\), \(\mathcal R_X^*=\mathcal R_X\), and the sign from \(z\mapsto-z\) gives the required adjoints. The quadratic Hamiltonian of \(\ell_k\) is quantized with normal ordering. In Darboux coordinates this sends terms \(pp,pq,qq\), respectively, to \(\h\) times second derivatives, first-order vector fields, and multiplication by a quadratic polynomial divided by \(\h\). Apply the dilaton translation
\[
 q(z)=t(z)-z\,1_X.
\]
Up to the harmless simultaneous operator-sign convention, these are the positive operators \(L_k^X\) in Getzler's formulation \cite[Equation~(1.2)]{Getzler1999}. There is no central scalar correction for \(k>0\). We fix the signs by that reference when invoking the Virasoro commutators.

For clarity, the proof uses the following finite description rather than a coordinate expansion of every coefficient.

\begin{proposition}[Positive coefficient rule]\label{conv:positive-rule}
A labelled coefficient of \(L_k^XZ_X\), for fixed \(k>0\), is a finite linear combination of the following operations on the disconnected Gromov--Witten tensors:
\begin{enumerate}[label=\textup{(\roman*)}]
\item add one special unit marking, with prescribed descendant and \(\mathcal R_X\)-powers;
\item act on one existing marking, changing its descendant exponent and applying an \(\mathcal R_X\)-power, with a coefficient polynomial in its first Hodge degree;
\item add two markings joined by the coevaluation of \((\, ,\,)_X\), with descendant powers, first-Hodge-degree polynomials, and \(\mathcal R_X\)-powers on its legs;
\item remove two primary markings, leaving their two cohomology inputs in the classical pairing with \(\mathcal R_X^{k+1}\), multiplied by the remaining disconnected Gromov--Witten tensor.
\end{enumerate}
The numerical coefficients and genus powers are universal for targets of fixed dimension. At most two existing labels are used nonordinarily in one summand.
\end{proposition}

\begin{proof}
Expand \(\ell_0(z\ell_0)^k\) using
\([z\partial_z,z]=z\) and \([\mu_X,\mathcal R_X]=\mathcal R_X\).
The resulting coefficients are polynomials in the first Hodge index, composed with Chern-class powers. The three blocks of its quadratic Hamiltonian give the last three types of operation under normal ordering; the dilaton shift gives the special unit marking. The block with both variables on the positive polarization has its only relevant negative shift at the primary pair, giving the classical \(\mathcal R_X^{k+1}\)-pairing. This is also the three-block decomposition in the cited coordinate formula. Since the Hamiltonian is quadratic, at most two existing labels are involved. Super quantization uses the same calculation with graded derivatives and permutations.
\end{proof}

The two lower operators are the standard string operator \(L_{-1}\) and the quantization of \(\ell_0\) with scalar
\[
 C_X=\frac{\chi(X)}{16}-\frac14\str(\mu_X^2),
\]
in Getzler's first-Hodge convention; its Chern-number form appears in \eqref{ind:central-normalization}. The nonpositive constraints require no new transport theorem. The string equation gives \(L_{-1}Z_X=0\). Together with \(L_1Z_X=0\), the commutator \([L_1,L_{-1}]=2L_0\) gives \(L_0Z_X=0\) with this normalization. Equivalently one may use the usual dilaton, homogeneity, divisor, and genus-one Riemann--Roch identities. We therefore prove all positive constraints.

\subsection{Admissible scalar tests}

An \emph{ordinary replica} means a separate copy of a disconnected Gromov--Witten tensor, with its own genus and degree variables. A \emph{tested positive constraint} consists of one coefficient rule from \cref{conv:positive-rule}, applied to one specified replica, followed by a finite tensor contraction using:
\begin{enumerate}[label=\textup{(\alph*)}]
\item ordinary replicas and classical cup integrals;
\item single classes obtained from algebraic classes on the relevant total families, allowing complex linear combinations and including units and Chern-class powers, or from homogeneous cohomology classes of fixed smooth projective sources through algebraic correspondences extending over those total families and inducing flat Hodge maps on their smooth loci;
\item coevaluations joining arbitrary pairs of slots, possibly on different replicas;
\item polynomials, or functions on the finite spectrum, of the first Hodge degree on any exposed leg;
\item cups with admissible single classes and numerical coefficients.
\end{enumerate}
The arities of this test are fixed. Only the chosen replica is acted on by \(L_k\); the test is then applied linearly to its coefficients. The test may be different for each coefficient that is to be proved zero.

When auxiliary cap markings are later inserted in the chosen replica, \(L_k\) acts on those variables as well. Their ordinary occurrences are preintegrated into the evaluation correspondence. If the coefficient rule modifies one of them, or removes it into a classical pairing, it is exposed as an additional position. By \cref{conv:positive-rule}, there are at most two such positions. Auxiliary cap labels on other replicas remain ordinary.

All contractions can be decomposed into homogeneous Hodge types. For a single supplied by a fixed smooth projective source, keep its cohomology slot exposed until the extending algebraic correspondence has been specialized; the homogeneous input itself need not be algebraic. A topologically extending class alone is not sufficient: its Hodge components need not be flat. Nor do we assert tested transport for an arbitrary flat Hodge correspondence without the specified algebraic extension over the total family. The first-index operations are permitted only in balanced numerical diagrams: the types of the algebraic correspondences, cups, and pairings must match the total input and output types. The limiting-Hodge calculation below explains why these numerical contractions are constant in smooth families and admit the required graded specialization.

\begin{definition}[Coherent data]\label{conv:coherent}
Admissible singles and numerical divisor-degree tests on several degenerations are \emph{coherent} if their restrictions agree on the common local component and stratum diagrams used in the comparison. For source-supplied singles this agreement is required for the specified extending algebraic correspondences with their cohomology slots exposed, before applying the fixed inputs. A class supported at a specified unaffected end is taken to have zero restriction on the other ends and replacement caps.
\end{definition}

The Chern-class insertion always has a coherent choice. The extension
\(-c_1(\omega_{\cY/B})\) restricts on a component \(Y_i\) with boundary \(D_i\) to \(c_1(TY_i)-[D_i]\), the logarithmic first Chern class. The local calculations use this class. They do not replace it silently by \(c_1(TY_i)\).

The tensor manipulations do not require individual primitive insertions to extend. They use full coevaluations and finitely many coherent singles; the final detection argument recovers arbitrary primitive inputs only on the target being proved.
